\documentclass[aspectratio=169]{beamer}

% Shared Beamer settings for Elements of AIML lectures (UPES).
% Keep this file free of \documentclass to allow reuse via \input.

% Allow lecture sources to use shared theme/assets without copying them into each lecture folder.
% NOTE: These paths are relative to `Unit-xx/Lecture-yy/latex/` (and `Lab/Experiment-xx/latex/`).
\makeatletter
\def\input@path{{../../../_shared/}{../../../_shared/UPES-Beamer-Theme/}}
\makeatother

\usetheme{UPES}
\setbeamertemplate{navigation symbols}{}

\usepackage{amsmath}
\usepackage{amssymb}
\usepackage{booktabs}
\usepackage{graphicx}
\usepackage{hyperref}
\usepackage{xcolor}
\usepackage{listings}

% Code listings (general-purpose).
\lstset{
  basicstyle=\ttfamily\small,
  keywordstyle=\color{blue},
  commentstyle=\color{gray},
  stringstyle=\color{teal},
  showstringspaces=false,
  columns=fullflexible,
  breaklines=true
}

\usepackage{tikz}
\usetikzlibrary{arrows.meta,positioning,shapes.geometric,calc}

% Per-slide source line (references shown on the slide itself).
\newcommand{\slidesources}[1]{%
  \vfill
  {\tiny\color{gray}\raggedright\sloppy\parbox{\linewidth}{Sources: #1}\par}%
}

% Unit-wise source strings for this subject (used by slides/notes).
% Path is relative to the lecture `latex/` folder (the compilation CWD).
% Source strings for Elements of AIML (used by slides + notes).

\newcommand{\AIMLRepoURL}{https://github.com/tali7c/Elements-of-AIML}

\newcommand{\AIMLLabSources}{%
Provost and Fawcett, \textit{Data Science for Business}.;%
McKinney, \textit{Python for Data Analysis}.;%
WEKA Documentation (Waikato Environment for Knowledge Analysis), accessed 2026-02-28.;%
Matplotlib and Seaborn documentation, accessed 2026-02-28.%
}

\newcommand{\AIMLUnitISources}{%
Rich and Knight, \textit{Artificial Intelligence}, McGraw Hill, 2017.;%
Russell and Norvig, \textit{Artificial Intelligence: A Modern Approach} (4e), Ch. 1--2 (Introduction, Intelligent Agents).;%
Poole and Mackworth, \textit{Artificial Intelligence: Foundations of Computational Agents} (2e), selected sections.%
}

\newcommand{\AIMLUnitIISources}{%
Russell and Norvig, \textit{Artificial Intelligence: A Modern Approach} (4e), Ch. 7--9 (Logical Agents, First-Order Logic, Inference).;%
Huth and Ryan, \textit{Logic in Computer Science} (2e), selected sections on propositional and first-order logic.;%
Genesereth and Nilsson, \textit{Logical Foundations of Artificial Intelligence}, selected chapters.%
}

\newcommand{\AIMLUnitIIISources}{%
Mueller and Massaron, \textit{Machine Learning for Dummies}, For Dummies, 2016.;%
Mitchell, \textit{Machine Learning}, selected chapters on learning basics and evaluation.;%
Scikit-learn documentation, accessed 2026-02-28.%
}

\newcommand{\AIMLUnitIVSources}{%
Mueller and Massaron, \textit{Machine Learning for Dummies}, For Dummies, 2016.;%
James, Witten, Hastie, and Tibshirani, \textit{An Introduction to Statistical Learning} (2e), selected chapters.;%
Scikit-learn documentation, accessed 2026-02-28.%
}

\newcommand{\AIMLUnitVSources}{%
NIST, \textit{AI Risk Management Framework (AI RMF 1.0)}, 2023.;%
Barocas, Hardt, and Narayanan, \textit{Fairness and Machine Learning} (online book), selected sections.;%
Knaflic, \textit{Storytelling with Data}, selected chapters on communicating results.%
}



\title[Elements of AIML]{Elements of AIML}
\subtitle{Unit 04 -- Lecture 01: Supervised Learning (Regression)}
\author{Tofik Ali}
\institute{School of Computer Science, UPES Dehradun}
\date{\today}

\graphicspath{{../images/}{../../../_shared/}{../../../_shared/UPES-Beamer-Theme/}}

\begin{document}

\UPESCoverFrame
\setcounter{framenumber}{0}

\begin{frame}
  \titlepage
  \vspace{0.3em}
  \begin{center}
    \footnotesize Topic focus: predicting \textbf{continuous values} and evaluating regression models\\[0.25em]
    \footnotesize Repository: \texttt{\AIMLRepoURL}
  \end{center}
  \slidesources{\AIMLUnitIVSources}
\end{frame}

\begin{frame}{Quick Links}
  \centering
  \hyperlink{sec:reg}{\beamerbutton{Regression}}\hspace{0.6em}
  \hyperlink{sec:linear}{\beamerbutton{Linear Model}}\hspace{0.6em}
  \hyperlink{sec:metrics}{\beamerbutton{Metrics}}\hspace{0.6em}
  \hyperlink{sec:overfit}{\beamerbutton{Overfitting}}\hspace{0.6em}
  \hyperlink{sec:summary}{\beamerbutton{Summary}}
\end{frame}

\begin{frame}{Agenda}
  \tableofcontents
\end{frame}

\section{Regression basics}
\label{sec:reg}

\begin{frame}{Learning Outcomes}
  By the end of this lecture, you should be able to:
  \begin{itemize}[<+->]
    \item Define regression and give common use-cases.
    \item Describe linear regression and the idea of a loss function.
    \item Compare MAE, MSE/RMSE, and $R^2$.
    \item Explain overfitting and how regularization helps.
  \end{itemize}
\end{frame}

\begin{frame}{Abbreviations \& Symbols}
  \begin{itemize}[<+->]
    \item $X$: features,\; $y$: target,\; $\hat{y}$: prediction
    \item $n$: samples,\; $d$: features,\; $w$: weights,\; $w_0$: intercept
    \item MAE/MSE/RMSE: regression error metrics
    \item $R^2$: coefficient of determination (variance explained)
    \item $\lambda$: regularization strength; $\|w\|_1,\|w\|_2$: L1/L2 penalties
  \end{itemize}
\end{frame}

\begin{frame}{What is Regression?}
  \begin{block}{Regression}
    Predict a \textbf{continuous} target value from input features.
  \end{block}
  \vspace{0.4em}
  Examples:
  \begin{itemize}[<+->]
    \item house price prediction
    \item demand forecasting
    \item temperature prediction
  \end{itemize}
\end{frame}

\section{Linear regression}
\label{sec:linear}

\begin{frame}{Linear Regression Model}
  \[
    \hat{y} = w_0 + w_1 x_1 + \cdots + w_d x_d
  \]
  \begin{itemize}[<+->]
    \item $w$ are learned from training data.
    \item The model is ``linear in parameters''.
  \end{itemize}
\end{frame}

\begin{frame}{Loss Function (MSE)}
  We learn parameters by minimizing a loss.
  \[
    \text{MSE} = \frac{1}{n}\sum_{i=1}^{n}(y_i - \hat{y}_i)^2
  \]
  \begin{itemize}[<+->]
    \item Squared error penalizes large mistakes more.
    \item Other losses exist (MAE, Huber), depending on needs.
  \end{itemize}
\end{frame}

\begin{frame}{Mini Worked Example (Compute Errors)}
  Suppose true values $y=[50,60,80]$ and predictions $\hat{y}=[55,58,70]$.
  \begin{itemize}[<+->]
    \item Errors: $e = y-\hat{y} = [-5,2,10]$
    \item MAE $= \frac{|{-5}|+|2|+|10|}{3}=\frac{17}{3}\approx 5.67$
    \item MSE $= \frac{25+4+100}{3}=\frac{129}{3}=43$
    \item RMSE $= \sqrt{43}\approx 6.56$ (same units as $y$)
  \end{itemize}
\end{frame}

\section{Metrics}
\label{sec:metrics}

\begin{frame}{Regression Evaluation Metrics}
  \begin{itemize}[<+->]
    \item \textbf{MAE}: average absolute error
    \item \textbf{MSE / RMSE}: squared error (RMSE is in original units)
    \item \textbf{$R^2$}: proportion of variance explained (can be misleading)
  \end{itemize}
\end{frame}

\section{Overfitting}
\label{sec:overfit}

\begin{frame}{Underfitting vs Overfitting}
  \begin{itemize}[<+->]
    \item \textbf{Underfitting}: model too simple, high bias, poor train and test.
    \item \textbf{Overfitting}: model too complex, low train error but high test error.
  \end{itemize}
  \vspace{0.4em}
  Remedies: more data, better features, regularization, cross-validation.
\end{frame}

\begin{frame}{Regularization (Idea)}
  Add a penalty to discourage overly large weights:
  \[
    \text{Loss} = \text{MSE} + \lambda \|w\|^2
  \]
  \begin{itemize}[<+->]
    \item Ridge ($\ell_2$) shrinks weights.
    \item Lasso ($\ell_1$) can push some weights to zero (feature selection effect).
  \end{itemize}
\end{frame}

\section{Summary}
\label{sec:summary}

\begin{frame}{Summary}
  \begin{itemize}
    \item Regression predicts continuous values.
    \item Linear regression minimizes a loss such as MSE.
    \item Metrics: MAE, RMSE, $R^2$ (interpret carefully).
    \item Overfitting can be reduced with regularization and proper validation.
  \end{itemize}
  \slidesources{\AIMLUnitIVSources}
\end{frame}

\begin{frame}{Exit Questions}
  \begin{enumerate}
    \item What is the difference between regression and classification?
    \item Why is RMSE often easier to interpret than MSE?
    \item What does $R^2$ measure and when can it be misleading?
    \item Explain overfitting and list two mitigation strategies.
  \end{enumerate}
\end{frame}

\end{document}
